% Generated draft: requires mathematical and source review.
\documentclass[11pt]{article}
\usepackage[margin=25mm]{geometry}
\usepackage{fontspec}
\setmainfont{Latin Modern Roman}
\usepackage{amsmath,amssymb,mathtools,unicode-math}
\setmathfont{Latin Modern Math}
\usepackage{enumitem,needspace,xurl,hyperref}
\setlength{\parindent}{0pt}
\setlength{\parskip}{5pt}
\setlength{\emergencystretch}{3em}
\newcommand{\N}{\mathbb{N}}
\newcommand{\Z}{\mathbb{Z}}
\newcommand{\Q}{\mathbb{Q}}
\newcommand{\R}{\mathbb{R}}
\newcommand{\C}{\mathbb{C}}
\newcommand{\F}{\mathbb{F}}
\newcommand{\A}{\mathbb{A}}
\newcommand{\PP}{\mathbb{P}}
\newcommand{\E}{\mathbb{E}}
\newcommand{\eps}{\varepsilon}
\newcommand{\dd}{\,\mathrm{d}}
\begin{document}
\section*{UnsolvedMath ID: 700996}
\textbf{Source problem code:} catalogue entry
\subsection{Definitions and mathematical statement}
Let \(R=\Q[x,y]\) be the polynomial ring in two commuting indeterminates over \(\Q\). The supplied source has an unmatched brace in its expression for this ring; the notation here reads that expression as \(\Q[x,y]\). Define
\[
G=\operatorname{SL}_2(R)=\{M\in M_2(R):\det M=1\},
\]
and define the elementary subgroup by
\[
E_2(R)=\left\langle
\begin{pmatrix}1&r\\0&1\end{pmatrix},
\begin{pmatrix}1&0\\r&1\end{pmatrix}
\;:\;r\in R
\right\rangle.
\]
For a subset \(A\subseteq G\), the notation \(\langle E_2(R),A\rangle\) denotes the smallest subgroup of \(G\) containing \(E_2(R)\cup A\). The source states as background that \(E_2(R)\ne G\).

\textbf{Problem (Kourovka Notebook 14.11).} Find a minimal subset \(A\subseteq G\) such that
\[
\langle E_2(R),A\rangle=G.
\]
The source does not specify the meaning of ``minimal.'' Inclusion-minimality would require
\[
\langle E_2(R),B\rangle\ne G\qquad\text{for every }B\subsetneq A.
\]
Minimum cardinality would instead require
\[
|A|\le |B|\qquad\text{for every }B\subseteq G\text{ with }\langle E_2(R),B\rangle=G.
\]
These are distinct interpretations; the supplied record does not determine which is intended. No finiteness requirement on \(A\) is stated.
\subsection{Short English statement}
Find a minimal collection of determinant-one matrices over \(\Q[x,y]\) which, together with all elementary matrices, generates \(\operatorname{SL}_2(\Q[x,y])\). The source leaves unclear whether minimality means inclusion-minimality or minimum cardinality.
\subsection{Sources}
\begin{itemize}
\item Local UnsolvedMath registry, permanent ID \(700996\), exact source code \texttt{catalogue entry}, title \emph{Kourovka Notebook Problem 14.11}; supplied upstream revision \texttt{d371ee917ac9647fbb06cd6fc529deec6f3dda82}.
\item Supplied source attribution: \emph{The Kourovka Notebook (Unsolved Problems in Group Theory)}, 14th issue (1999), Problem 14.11; also identified in the record with the 21st issue edition (2026), arXiv identifier \texttt{1401.0300v48}. The collection description names E. I. Khukhro and V. D. Mazurov as editors. These bibliographic details are supplied metadata and have not been independently verified.
\item The background and proposer field credit V. G. Bardakov, with proposed year 1999. The statement itself begins with the parenthetical attribution ``Yu. I. Merzlyakov.'' Both attributions are retained; their respective roles are not explained in the supplied record.
\item The supplied registry labels the problem \emph{open}. This is a source status label, not an independently verified assessment of the literature.
\end{itemize}
\end{document}
